% Template:     Informe LaTeX
% Documento:    Engineering report - Static load balancing for a CPU+GPU stencil
% Codificación: UTF-8
%
% Standalone engineering report on the HybridEngine (CPU+GPU, Divisible Load
% Theory, Barlas 11.3) sweep results. Build like the other two documents:
%   cd report && latexmk -pdf hybrid_report.tex
% Replaces the earlier Markdown note hybrid_load_balancing_analysis.md.

% CREACIÓN DEL DOCUMENTO
\documentclass[
	english,
	letterpaper, oneside
]{article}

% INFORMACIÓN DEL DOCUMENTO
\def\documenttitle {Static Load Balancing for a CPU\,+\,GPU Stencil}
\def\documentsubtitle {Static load balancing across CPU, iGPU and dGPU: a verdict}
\def\documentsubject {Computación en GPU}

\def\documentauthor {Matías Saavedra}
\def\coursename {Computación en GPU}
\def\coursecode {CC7515-1}

\def\universityname {Universidad de Chile}
\def\universityfaculty {Facultad de Ciencias Físicas y Matemáticas}
\def\universitydepartment {Departamento de Ciencias de la Computación}
\def\universitydepartmentimage {departamentos/dcc}
\def\universitydepartmentimagecfg {height=1.57cm}
\def\universitylocation {Santiago de Chile}

% INTEGRANTES, PROFESORES Y FECHAS
\def\authortable {
	\begin{tabular}{ll}
		Author:
		& \begin{tabular}[t]{l}
			Matías Saavedra
		\end{tabular} \\
		& \\
		\multicolumn{2}{l}{Document: engineering report (hybrid load balancing)} \\
		\multicolumn{2}{l}{Last updated: \today} \\
		\multicolumn{2}{l}{\universitylocation}
	\end{tabular}
}

% IMPORTACIÓN DEL TEMPLATE
\input{template}

% INICIO DE PÁGINAS
\begin{document}

% PORTADA
\templatePortrait

% CONFIGURACIÓN DE PÁGINA Y ENCABEZADOS
\templatePagecfg

% RESUMEN O ABSTRACT
\begin{abstractd}
	The \texttt{HybridEngine} runs several compute nodes on one Game of Life board
	(HighLife, \texttt{B36/S23}), split by rows with Divisible Load Theory (DLT, Barlas
	\S11.3). This report analyses static load balancing on a GTX~1660~Ti Max-Q, first
	paired with a 12-thread CPU and then with the machine's Intel UHD~630 integrated GPU.
	The CPU+GPU pair is a textbook ``one node dominates'' regime: the GPU is
	$\sim$160$\times$ the CPU, so the DLT-optimal CPU share is 0.6\,\% and the balancing
	ceiling is smaller than the per-step coordination cost --- the hybrid does not beat
	pure GPU, and a raw sweep that appears to is a kernel confound (the hybrid runs a
	faster, branch-free halo kernel). The apparent ``OpenCL regression'' seen during that
	run was a misdiagnosis: the slow 1.7\,Gcells/s figure is the integrated GPU, which the
	default OpenCL selection picks as platform~0, not a driver fault. Reusing that idle
	iGPU as a genuine second DLT node changes the verdict. Because the inter-node seam is
	a single fixed row while compute grows as $N^2$, the communication cost vanishes as
	$1/N$: a well-chosen static iGPU\,$+$\,dGPU split beats pure dGPU by a margin that
	grows with the grid, reaching $+5.5\,\%$ at $N=32768$ and 99\,\% of the
	$R_\text{iGPU}+R_\text{dGPU}$ ceiling. The prize predicted by DLT is real once the
	load is large enough to amortize the seam --- though the default auto-calibration
	still overshoots the optimum.
\end{abstractd}

% TABLA DE CONTENIDOS - ÍNDICE
\templateIndex

% CONFIGURACIONES FINALES
\templateFinalcfg

% ======================= INICIO DEL DOCUMENTO =======================

\section{Scope and setup}
\label{sec:setup}

Backend under test: the \texttt{HybridEngine} (branch
\texttt{static\_load\_balancing}) --- CPU\,+\,GPU on one board, split by rows, the
partition chosen by DLT. Hardware: an NVIDIA GTX~1660~Ti Max-Q (\texttt{sm\_75},
60\,W power ceiling) and a 12-thread CPU, on Linux. Data:
\texttt{results/linux/sweep\_\{gpu\_opt,scaling\}.csv} from \texttt{scripts/sweep.sh};
figure \texttt{img/hybrid\_load\_balancing.png}.

\begin{itemize}
	\item Metric. Cells evaluated per second,
	$\text{Mcells/s} = (\text{rows}\times\text{cols}\times\text{gens}) / (t\times10^{6})$,
	measured with the \texttt{NullRenderer} so host\,$\leftrightarrow$\,device
	transfers and drawing never enter the timed loop. Timing is in-program and
	uniform per backend (CPU and hybrid: \texttt{std::chrono}; CUDA:
	\texttt{cudaEvent}). The hybrid figure is the wall of the concurrent CPU+GPU
	region, i.e.\ $\approx$ the slower of the two nodes.
	\item Sweep. Experiment A is the GPU block\,$\times$\,shared/local sweep;
	experiment B is the scaling vs.\ $N$ for CPU(threads) / CUDA / OpenCL / hybrid.
	We report the best (highest-Mcells) row per configuration, matching ``mejor
	tiempo'' in the graded report. Bounded edges, deterministic random fill
	(\texttt{--seed 1}: identical work across backends).
	\item Controlled follow-ups (this analysis). Single-knob runs that isolate why
	the hybrid behaves as it does: pure CUDA vs.\ the halo kernel on the full board
	(\texttt{--cpu-frac 0}) vs.\ the auto-split hybrid, all at block 128 so the only
	variable is the code path.
\end{itemize}

GPU health was verified before and after the sweep (idle 4.9\,W, 39--41\,\textdegree
C, 60\,W ceiling, no power-cap creep, no throttle flags), so the GPU numbers are not
thermally compromised. Run-to-run boost-clock variance is real, though --- see
\S\ref{sec:variance}.

\section{Throughput results}
\label{sec:throughput}

\begin{table}[H]
	\centering
	\caption{Best throughput per backend per grid size $N$ (Gcells/s, kernel time).
	The last column, ``GPU halo-only'', is the hybrid's GPU path run with
	\texttt{--cpu-frac 0} (halo kernel, full board, no CPU) --- a control that
	isolates the kernel from the load balancing.}
	\label{tab:throughput}
	\begin{tabular}{@{}rrrrrr@{}}
		\toprule
		$N$ & CPU (best $t$) & CUDA & OpenCL & Hybrid (auto) & GPU halo-only \\
		\midrule
		512   & 0.151 & 24.76 & 1.67 & 5.20  & 14.57 \\
		1024  & 0.146 & 27.55 & 1.70 & 9.56  & 25.07 \\
		2048  & 0.144 & 31.52 & 1.70 & 19.91 & 30.70 \\
		4096  & 0.142 & 32.25 & 1.69 & 22.75 & 31.42 \\
		8192  & 0.147 & 26.20 & 1.69 & 25.12 & 28.79 \\
		16384 & 0.152 & 24.83 & 1.68 & 28.46 & 28.72 \\
		\bottomrule
	\end{tabular}
\end{table}

The headlines from the table:

\begin{itemize}
	\item The CPU plateaus at $\sim$0.15\,Gcells/s regardless of $N$: a 1-byte
	9-point stencil is memory-bandwidth bound, and beyond $\sim$6 threads the cores
	contend for the same bandwidth (scaling is flat, not linear). This sets
	$R_\mathrm{cpu}$.
	\item CUDA is $\sim$160--230$\times$ the CPU, peaking $\sim$32\,Gcells/s at
	$N\approx2048$--$4096$, then easing to $\sim$25 at $N=16384$ (larger working sets
	spill L2, so the kernel becomes more memory-bound).
	\item The ``OpenCL'' column reads 1.7\,Gcells/s, $\sim$15$\times$ below CUDA. This is
	not a slow NVIDIA path but the Intel UHD~630 integrated GPU, which the default OpenCL
	selection picks as platform~0 (see \S\ref{sec:opencl}); pinned to NVIDIA, OpenCL
	matches CUDA. That idle iGPU becomes the second node of the two-GPU study
	(\S\ref{sec:twogpu}).
\end{itemize}

\section{The load-balancing split (DLT \S11.3)}
\label{sec:split}

The calibration phase times a few CPU-only and GPU-only steps on the actual grid
and applies the all-finish-together optimum:

\begin{equation}
	\label{eq:dlt}
	\text{part}_\mathrm{gpu} = \frac{R_\mathrm{gpu}}{R_\mathrm{cpu} + R_\mathrm{gpu}},
	\qquad
	\text{part}_\mathrm{cpu} = 1 - \text{part}_\mathrm{gpu},
\end{equation}

\noindent with $R = 1/p$ the measured throughput (cells/s), converted to a row split
$s = \text{round}(\text{part}_\mathrm{cpu}\cdot R)$.

\begin{table}[H]
	\centering
	\caption{Calibrated rates and the resulting split. $R$ in Mcells/s. The split is
	stable at 0.6\,\% CPU at every size, exactly as $R_\mathrm{cpu}/R_\mathrm{gpu}
	\approx 0.006$ predicts.}
	\label{tab:split}
	\begin{tabular}{@{}rrrrrrr@{}}
		\toprule
		$N$ & $R_\mathrm{cpu}$ & $R_\mathrm{gpu}$ & GPU/CPU & CPU rows & CPU share & DLT bound $R_\mathrm{cpu}{+}R_\mathrm{gpu}$ \\
		\midrule
		2048  & 157.6 & 24\,504 & 156$\times$ & 13 & 0.6\,\% & 24\,661 \\
		4096  & 163.8 & 27\,176 & 166$\times$ & 25 & 0.6\,\% & 27\,340 \\
		8192  & 158.5 & 25\,271 & 159$\times$ & 51 & 0.6\,\% & 25\,430 \\
		16384 & 158.0 & 27\,286 & 173$\times$ & 94 & 0.6\,\% & 27\,444 \\
		\bottomrule
	\end{tabular}
\end{table}

The split matches the \S11.3 optimum to the row, and the two nodes finish together:
at $N=16384$ the CPU needs $94 \times 16384 / 158$\,Mcells/s $\approx 9.74$\,ms and
the GPU needs $16\,290 \times 16384 / 27\,286$\,Mcells/s $\approx 9.78$\,ms. The DLT
machinery works precisely as the book describes --- this is not a tuning failure.

The catch is the ceiling. Because the CPU's optimal share is only 0.6\,\%, the DLT
``equivalent node'' throughput $R_\mathrm{cpu}+R_\mathrm{gpu}$ is at most 0.6\,\%
above the GPU alone. That is the entire prize, and it is tiny.

\section{The figure}
\label{sec:figure}

\insertimage{hybrid_load_balancing}{width=\linewidth}{Hybrid scaling. Left (log--log,
all backends): the GPU backends sit $\sim$160--230$\times$ above the flat CPU line;
the OpenCL curve at $\sim$1.7\,Gcells/s is the Intel iGPU (the default OpenCL device),
not a slow NVIDIA path. Right (GPU band): CUDA
(\texttt{life\_global}) peaks early then sags; the halo kernel holds
$\sim$29\,Gcells/s at large $N$; the auto-split hybrid climbs from far below toward
the GPU band as the per-step overhead amortizes, only catching up near $N=16384$.}

\section{Unexpected findings}
\label{sec:findings}

\subsection{``Hybrid beats CUDA at large N'' is a kernel effect, not load balancing}
\label{sec:confound}

The raw sweep shows the hybrid mean exceeding the CUDA mean at $N=16384$ (27.4 vs.\
24.6\,Gcells/s, $+11$\,\%) with low variance --- too consistent to be noise, yet
impossible to attribute to the CPU, whose DLT-optimal contribution is only 0.6\,\%. A
controlled run resolves it (block 128, full board):

\begin{table}[H]
	\centering
	\caption{Like-for-like at block 128, mean of repeated runs (Gcells/s). The halo
	kernel beats \texttt{life\_global} with zero CPU; adding the CPU and the
	coordination (hybrid auto) gives it all back.}
	\label{tab:kernel}
	\begin{tabular}{@{}rrrr@{}}
		\toprule
		$N$ & CUDA \texttt{life\_global} & Halo kernel (\texttt{--cpu-frac 0}, no CPU) & Hybrid (auto) \\
		\midrule
		4096  & 25.68 & 26.61 ($+3.6$\,\%) & 18.25 ($-29$\,\%) \\
		16384 & 25.14 & 28.12 ($+12$\,\%)  & 25.39 ($+1$\,\%) \\
		\bottomrule
	\end{tabular}
\end{table}

The halo kernel is intrinsically faster than \texttt{life\_global}, with zero CPU
involvement. The reason: \texttt{life\_global} does a per-cell vertical bounds check
(\texttt{if (nr < 0 || nr >= rows)}, or a modulo under \texttt{--wrap}) for every
neighbour row; the halo kernel reads its vertical neighbours unconditionally from
always-present ghost rows, so it is branch-free in the vertical direction. At large
$N$ the kernel is compute/branch-bound (not launch-bound), so removing those branches
buys 10--16\,\%. The ``hybrid beats CUDA'' headline was the hybrid's GPU slice quietly
running a better kernel --- a confound, not a balancing win.

\subsection{Once the kernel is held constant, load balancing is a net loss}
\label{sec:netloss}

Comparing like-for-like (both on the halo kernel), hybrid-auto is below halo-only at
every size --- $-29$\,\% at $N=4096$, $\approx$ break-even at $N=16384$. The per-step
coordination (ghost-row exchange $+$ waking the worker pool $+$ synchronising the two
nodes, $\sim$0.3--1\,ms/step) costs more than the 0.6\,\% the CPU offloads. Concretely
at $N=16384$: the DLT-optimal split saves the GPU $\sim$0.6\,\% of its time
($\sim$0.06\,ms) by handing 94 rows to the CPU, while the coordination adds
$\sim$0.6--1\,ms --- a clear net loss that only looks like break-even because the halo
kernel's own speed-up cancels it.

\subsection{The apparent ``OpenCL regression'' was the integrated GPU}
\label{sec:opencl}

This box previously measured OpenCL at $\sim$37\,Gcells/s, on par with CUDA (both
lower through the same PTX/SASS on NVIDIA). A later sweep measured 1.7\,Gcells/s and
was first recorded as a driver/ICD regression. It was not. With explicit device
pinning (the \texttt{GOL\_OCL\_DEVICE} environment variable, added in
\texttt{OclDeviceSelect.hpp}), the 1.7\,Gcells/s figure is the Intel UHD~630
integrated GPU, which the default device selection picks because it enumerates as
OpenCL platform~0. Pinned to NVIDIA, OpenCL runs at $\sim$23--28\,Gcells/s, matching
CUDA. So there was no regression and no reboot is needed --- only a silent device
choice. This reframes the exercise: the ``slow OpenCL'' was a second, real GPU
sitting idle, which motivates the two-GPU split of \S\ref{sec:twogpu}.

\subsection{GPU run-to-run variance is $\sim$10--25\,\%}
\label{sec:variance}

A single CUDA run at $N=4096$ gave 24.9\,Gcells/s vs.\ the best-of-5 of 31.8 ---
boost-clock and scheduling variance, largest at small $N$ (per-rep CV up to $\sim$30\,\%
at $N=2048$, down to $\sim$1--5\,\% at $N=16384$). Any margin smaller than this band ---
including the entire 0.6\,\% load-balancing prize --- is unmeasurable. Report best-of-$N$
and treat sub-10\,\% differences at small grids as noise.

\subsection{OpenCL work-group limit at block 512}
\label{sec:wg512}

\texttt{--block 512} fails OpenCL with \texttt{CL\_INVALID\_WORK\_GROUP\_SIZE} (code
$-54$) on this device (CUDA runs 512 fine). The sweep skips it gracefully; it is a
device limit, not a regression. Block 64--128 is optimal on both engines anyway.

\section{From CPU+GPU to iGPU+dGPU}
\label{sec:twogpu}

The negligible-CPU verdict says the CPU is the wrong second node; the corrected
OpenCL finding (\S\ref{sec:opencl}) supplies a better one --- the Intel UHD~630
integrated GPU, reachable only through OpenCL. The engine was generalized from a
fixed [CPU, GPU] pair to an ordered list of nodes, each a CPU pool, a CUDA
discrete-GPU partition, or an OpenCL partition pinned to a chosen device. The default
composition is iGPU\,$+$\,dGPU, with the CPU off by default; adjacent nodes exchange
the same one-row ghost halo, and inter-GPU seams stage through host memory (there is no
cross-vendor peer copy). Every composition (\texttt{igpu,dgpu}, \texttt{cpu,dgpu}, the
three-way \texttt{cpu,igpu,dgpu}) verifies bit-for-bit against the CPU oracle in both
edge modes.

\subsection{The DLT model, applied}
\label{sec:twogpu-model}

In the terms of \S11.3, the board is the divisible load $L = R\cdot C$ cells, node $i$
receives a row share $\mathit{part}_i$, and its inverse speed is $p_i = 1/R_i$. The
book's affine node cost has a computation term $p_i(\mathit{part}_i L + e_i)$ (Eq.~11.5)
and a distribution term $l_i(a\,\mathit{part}_i L + b)$ (Eq.~11.7). Our ghost seam is
exactly the constant $b$ of that model: a single boundary row moved per step,
independent of $\mathit{part}_i$ --- the ``boundary row of pixels'' of the book's
Fig.~11.5(a). The optimum (Fig.~11.9) makes both nodes finish together.

Dropping the communication constant $b$ and the fixed overhead $e_i$ leaves the naive
closed form, which is what the auto-calibration computes:
\begin{equation}
	\label{eq:naive2gpu}
	\mathit{part}_\text{iGPU} = \frac{R_\text{iGPU}}{R_\text{iGPU}+R_\text{dGPU}}
	\approx \frac{2.1}{2.1+31.5} \approx 6.3\,\%,
\end{equation}
with $R_\text{iGPU}\approx 2.1$ and $R_\text{dGPU}\approx 31.5$\,Gcells/s (the
pure-dGPU halo rate). This is an order of magnitude larger than the CPU's 0.6\,\%
share, so the DLT ceiling $R_\text{iGPU}+R_\text{dGPU}\approx 33.6$\,Gcells/s now sits
$\sim$6.7\,\% above the discrete GPU alone --- a prize worth chasing.

\subsection{Theory versus the empirical optimum}
\label{sec:twogpu-sweep}

Fig.~\ref{fig:frac} sweeps the iGPU row share directly, holding the composition at
\texttt{igpu,dgpu} and measuring steady-state throughput (best of five, explicit
\texttt{--fracs} so no calibration bias, long runs so the boost clocks are ramped).
Each curve rises to a peak near a 5--6\,\% iGPU share and then falls off a cliff: past
$\sim$7\,\% the slow iGPU is handed more rows than it can finish inside the dGPU's step
time and becomes the bottleneck. The empirical peaks sit at 5--6\,\% (four of the five
grids peak at 6\,\%), straddling the naive 6.3\,\% of Eq.~\eqref{eq:naive2gpu} to within the sweep's 1\,\%
grid resolution: near the top the throughput curve is flat, so best-of-five boost
variance moves the winning grid point by $\pm$1 step. It is tempting to read the
sub-6.3\,\% peaks as the seam and the iGPU's launch overhead lowering its effective
rate --- but \S\ref{sec:twogpu-dlt} measures those terms and feeds them into the
reference DLT library, and the genuine correction to the optimum is under $0.05$
percentage points, far too small to explain a shift to 5\,\%. The sub-6.3\,\% readings
are sampling noise, not a real overhead-driven bias.

\begin{table}[H]
	\centering
	\caption{iGPU+dGPU fraction sweep (best of five, steady state). ``dGPU'' is the
	pure discrete-GPU halo rate ($\mathit{part}_\text{iGPU}=0$); ``best'' is the peak
	over the sweep at the listed iGPU share. Ceiling $=R_\text{iGPU}+R_\text{dGPU}$
	with $R_\text{iGPU}=2.1$. Rates in Gcells/s.}
	\label{tab:frac}
	\begin{tabular}{@{}rrrrrr@{}}
		\toprule
		$N$ & dGPU & best (iGPU share) & win & ceiling & \% of ceiling \\
		\midrule
		8192  & 30.43 & 30.77 (6\,\%) & $+1.1\,\%$ & 32.53 & 94.6\,\% \\
		16384 & 31.06 & 32.07 (6\,\%) & $+3.2\,\%$ & 33.16 & 96.7\,\% \\
		24576 & 31.58 & 32.73 (5\,\%) & $+3.6\,\%$ & 33.68 & 97.2\,\% \\
		32768 & 31.62 & 33.36 (6\,\%) & $+5.5\,\%$ & 33.72 & 98.9\,\% \\
		40960 & 31.54 & 33.17 (6\,\%) & $+5.2\,\%$ & 33.64 & 98.6\,\% \\
		\bottomrule
	\end{tabular}
\end{table}

\subsection{Why the win grows with the grid}
\label{sec:twogpu-scaling}

The decisive observation is the scaling (Table~\ref{tab:frac}, Fig.~\ref{fig:frac}(B)).
The seam $b$ is a single row --- a fixed $C$ bytes per step --- while the compute load
$L = R\cdot C$ grows as $N^2$. The communication-to-computation ratio is therefore
$\sim b/L \sim 1/N$ and vanishes as the grid grows, so the achievable throughput must
approach the DLT ceiling $R_\text{iGPU}+R_\text{dGPU}$. That is what happens: the best
static split beats pure dGPU by $+1.1\,\%$ at $N=8192$, growing monotonically to
$+5.5\,\%$ at $N=32768$ ($+5.2\,\%$ at $40960$, within variance of the peak). As a
fraction of the ceiling, the hybrid climbs from 94.6\,\% to 98.9\,\%, while pure dGPU
stays pinned at $\sim$93.7\,\% --- it is simply leaving the iGPU's $\sim$6\,\% on the
table. At the largest grids the two-GPU hybrid captures essentially the entire
theoretical bound: DLT behaving exactly as \S11.3 describes in the large-load limit,
where communication is negligible and the partition delivers the sum of the nodes'
rates.

\insertimage[\label{fig:frac}]{hybrid_frac_sweep}{width=\linewidth}{iGPU+dGPU fraction sweep.
(A)~Steady-state throughput vs the iGPU row share, one curve per grid size (best of
five; darker = larger $N$). Each curve rises to a peak at 5--6\,\% then falls off a
cliff as the iGPU is overloaded past the dGPU's step time --- steepest at the largest
grid; the peaks straddle the naive-DLT share $R_i/(R_i{+}R_d)\approx 6.3\,\%$ (dashed)
within the sweep's 1\,\% step. (B)~Best static split and pure dGPU as shares of the
$R_i+R_d$ ceiling, on one axis: the hybrid climbs from 94.6\,\% to 98.9\,\% while pure
dGPU stays pinned at $\sim$93.7\,\%; the gap between the two curves is the hybrid's
win over pure dGPU (annotated at the ends, $+1.1\,\%\to+5.5\,\%$), growing as the fixed
one-row seam is amortized against $N^2$ compute.}

This reverses the pessimistic reading from the CPU+GPU study and the small grids: the
earlier ``a few percent below pure dGPU'' was measured only at $N\le 16384$, where the
seam still costs an appreciable share and the margin hid inside the $\sim$5--10\,\%
boost variance. With clean measurement and large grids, a well-chosen static iGPU+dGPU
split is a real, if modest, win --- one whose size is governed by the divisible-load
communication term, not by luck.

\subsection{DLTlib cross-check: the overhead is negligible}
\label{sec:twogpu-dlt}

To separate genuine divisible-load effects from measurement noise, we fed the platform
into the reference library the engine is modelled on --- DLTlib (Barlas, Appendix~G) ---
and compared its optimum to ours. The platform maps to DLTlib's single-level tree as an
idle load-originating root (a pure distributor) with one compute child per device of
processing speed $\mathit{power}_i = 1/R_i$; the library's fixed-cost term $e_{0,i}$ is
in load units, so a physical per-step overhead $\tau_i$ (seconds) enters as
$e_{0,i} = \tau_i R_i$.

\paragraph{The naive split is the library's optimum.} With communication switched off
($\mathit{link}=0$, $e_0=0$) DLTlib's \texttt{SolveImage} returns
$\mathit{part}_i = R_i/\sum_j R_j$ to six digits --- the iGPU share $6.30\,\%$ and, for
the CPU+GPU pair, $0.58\,\%$. So the engine's hand-rolled Eq.~\eqref{eq:naive2gpu} is
not an approximation: it \emph{is} the \S11.3 optimum, and the library confirms it
exactly.

\paragraph{The measured overhead does not move it.} We then measured each device's
affine per-step timing $t = \tau_i + \text{cells}/R_i$ by running it solo across board
sizes (Fig.~\ref{fig:dltfit}): $R_\text{iGPU}=2.12$, $R_\text{dGPU}=31.5$\,Gcells/s
(matching the sweep), with fixed overheads of only $\tau_\text{iGPU}\approx23\,\mu$s and
$\tau_\text{dGPU}\approx6\,\mu$s. Feeding $(R_i,\tau_i)$ back into DLTlib gives the
overhead-aware optimum of Table~\ref{tab:dlt}: it stays within $0.05$ percentage points
of the naive $6.30\,\%$ at every grid size, and the correction vanishes as $1/N^2$.

\begin{table}[H]
	\centering
	\caption{DLTlib optimum with the measured fixed overhead ($\tau_\text{iGPU}=23\,\mu$s,
	$\tau_\text{dGPU}=6\,\mu$s) folded in via $e_0=\tau R$, versus the naive share and the
	fraction-sweep peak. The overhead correction is $\le 0.05$\,pp at every $N$.}
	\label{tab:dlt}
	\begin{tabular}{@{}rrrr@{}}
		\toprule
		$N$ & naive $R_i/\Sigma R$ & DLTlib $+$ fitted $e_0$ & empirical peak \\
		\midrule
		8192  & 6.30\,\% & 6.25\,\% & $\sim$6\,\% \\
		16384 & 6.30\,\% & 6.29\,\% & $\sim$6\,\% \\
		24576 & 6.30\,\% & 6.30\,\% & $\sim$5\,\% \\
		32768 & 6.30\,\% & 6.30\,\% & $\sim$6\,\% \\
		40960 & 6.30\,\% & 6.30\,\% & $\sim$6\,\% \\
		\bottomrule
	\end{tabular}
\end{table}

The reason is scale: at $N\ge 8192$ a step is milliseconds while the fixed overhead is
tens of microseconds, so $\tau/t_\text{step}\lesssim 10^{-2}$ and the split is
overhead-insensitive. The prediction is therefore a flat $6.3\,\%$, and the empirical
peaks straddle it within one sweep step (Fig.~\ref{fig:dltpred}). The reading that the
optimum ``sits below $6.3\,\%$ because of the seam and launch overhead'' does not survive
this measurement: neither the fixed launch cost nor the one-row seam is within two orders
of magnitude of what a $-1.3$\,pp shift would require. The $5$--$6\,\%$ scatter is grid
resolution and boost variance, not a physical bias. (As a consistency check, adding a
\emph{proportional} link cost instead moves the share the wrong way --- up, toward an
even split --- because DLTlib's link models transfer cost proportional to the load moved,
whereas the GoL seam is a fixed single row: the overhead belongs in $e_0$, not the link.)

\insertimage[\label{fig:dltfit}]{dlt_fit}{width=\linewidth}{Per-device affine timing
$t_\text{step}=\tau+\text{cells}/R$, each device run solo through the hybrid harness.
(A)~time per step vs.\ cells (log--log): the slope gives $R$. (B)~time per cell: the
fixed overhead $\tau$ shows only as the small-$N$ bump (leftmost point) and is gone by
$N=1024$, pinning $\tau_\text{iGPU}\approx23\,\mu$s, $\tau_\text{dGPU}\approx6\,\mu$s.}

\insertimage[\label{fig:dltpred}]{dlt_pred_vs_empirical}{width=0.85\linewidth}{DLTlib
optimum with fitted overhead (blue) vs.\ the fraction-sweep peaks (green, $\pm$1 sweep
step) and the naive $R_i/\Sigma R$ (dotted). The prediction is a flat $6.3\,\%$; the
empirical peaks straddle it within resolution --- confirming the split is the naive DLT
optimum and the sub-6.3\,\% readings are noise.}

\subsection{The catch: the auto-split overshoots}
\label{sec:twogpu-auto}

The win requires the correct fraction, and the default auto-calibration does not
reliably find it --- not because the naive Eq.~\eqref{eq:naive2gpu} is wrong (\S\ref{sec:twogpu-dlt}
shows it is the exact optimum) but because a short in-run calibration mis-measures the
rates that feed it. It is biased by boost warm-up: the dGPU needs a long ramp to reach
its 31.5\,Gcells/s peak (\S\ref{sec:twogpu-dlt}), so a brief calibration under-measures
it and inflates the iGPU share above the true 6.3\,\%. Because the throughput curve past
the optimum falls off a cliff while the approach is gentle, this asymmetry makes overshoot
far costlier than undershoot. Lengthening calibration (\texttt{--calib-steps 50}) helps
but does not fully cure it. So today the win is reachable with a hand-set
\texttt{--fracs 0.05,0.95}; making the default land there --- by calibrating the dGPU with
a proper boost warm-up so its rate is not under-measured, or simply by damping the split
toward the dGPU --- is the obvious next step.

\section{Conclusion: does static load balancing pay off here?}
\label{sec:conclusion}

The implementation is correct and the theory is vindicated. Whether it pays off
depends on the node pair and the grid size.

\begin{itemize}
	\item Correctness. The hybrid verifies bit-for-bit against the \texttt{CpuEngine}
	oracle in both edge modes and across all splits; the calibrated partition matches
	the \S11.3 optimum to the row; the two nodes finish together. As a teaching
	realization of DLT on a genuinely heterogeneous CPU+GPU pair (disjoint memory,
	real seam communication), it does exactly what the book says.
	\item But this is the textbook ``one node dominates'' regime. With the GPU
	$\sim$160$\times$ the CPU, the DLT-optimal CPU share is 0.6\,\%, so the
	theoretical ceiling over pure GPU is $\sim$0.6\,\% --- and the per-step ghost
	exchange $+$ pool wakeup $+$ sync exceed that. The hybrid ranges from $-29$\,\%
	(small/mid $N$, overhead dominated) to break-even (large $N$, overhead
	amortized). It never meaningfully beats pure GPU.
	\item One step up the device hierarchy, it does pay off. Replacing the CPU with the
	idle integrated GPU (\S\ref{sec:twogpu}) raises the minority share to $\sim$6\,\%,
	and because the seam is a single fixed row its cost vanishes as $1/N$. A well-chosen
	static iGPU+dGPU split then beats pure dGPU by a margin that grows with the grid ---
	up to $+5.5\,\%$ at $N=32768$, reaching 99\,\% of the $R_\text{iGPU}+R_\text{dGPU}$
	ceiling. The prize is real here; at small $N$ it hid inside the boost noise, and the
	default auto-split still overshoots it (\S\ref{sec:twogpu-auto}).
	\item The only speed-up in the data is a kernel-engineering side effect. The
	ghost-cell halo kernel, built to make domain decomposition possible, is
	incidentally 10--16\,\% faster than the original global kernel at large $N$
	because it trades vertical boundary branches for unconditional ghost-row reads.
\end{itemize}

\subsection*{Recommendations}

\begin{enumerate}
	\item Ship pure GPU with the halo (branch-free) kernel. That is the real lever
	this exercise surfaced --- a free 10--16\,\% at large $N$, no CPU, no
	coordination. Consider promoting the halo formulation to the default global
	kernel (it needs only a one-row zero-pad at the board's top/bottom under bounded
	edges).
	\item The hybrid pays off under either of two conditions: the node rates are within a
	small factor (a vectorized / bit-sliced CPU, a weaker / shared GPU), so the minority
	share is large; or the load is large enough that the fixed one-row seam is negligible
	against $N^2$ compute. The CPU+GPU pair meets neither (0.6\,\% share, and the grids do
	not go high enough); the iGPU\,$+$\,dGPU pair meets the second at $N\gtrsim 16384$ and
	wins by up to $+5.5\,\%$ (\S\ref{sec:twogpu}). Use it there with a hand-set
	\texttt{--fracs} near $0.05$--$0.06$.
	\item Fix the auto-calibrated split so the default \texttt{--engine hybrid} lands on
	the $\sim$5--6\,\% iGPU optimum instead of overshooting into the cliff: fold a
	communication/overhead term into the calibration, or damp the split toward the dGPU
	(\S\ref{sec:twogpu-auto}). There is no OpenCL pathology to fix --- the slow figure was
	the integrated GPU (\S\ref{sec:opencl}).
\end{enumerate}

The bottom line: static load balancing is the right tool for balanced heterogeneous
nodes, and the implementation proves it can be done correctly and cheaply. On a
GTX~1660~Ti paired with a scalar CPU stencil it has nothing to balance --- the honest
result is ``$\approx$ pure GPU'', and the interesting performance story is the kernel,
not the split.

\end{document}
